% §12 Exercises
%
% Cross-reference resolution status:
%
% DEFINED here:
%   \label{ex:ch12-1} through \label{ex:ch12-11}
%
% RESOLVED (used here, defined elsewhere):
%   \Cref{sec:sets-to-relations}                  -- §12.1 (Plan 12 Task 2)
%   \Cref{sec:relation-composition}               -- §12.2 (Plan 12 Task 3)
%   \Cref{thm:join-rates}                         -- §12.2 (Plan 12 Task 3)
%   \Cref{sec:approximate-joins}                  -- §12.3 (Plan 12 Task 4)
%   \Cref{prop:union-fpr-fnr}                     -- ch6 §6.2
%   \Cref{prop:intersection-fpr-fnr}              -- ch6 §6.2
%   \Cref{prop:minhash-collision}                 -- ch10 §10.4
%   \Cref{prop:minhash-error}                     -- ch10 §10.4
%   \Cref{def:minhash}                            -- ch10 §10.4
%   \Cref{sec:minhash-lsh}                        -- ch10 §10.4
%   \Cref{sec:entropy-lower-bound}                -- ch8 §8.3 (entropy floor)
%   \Cref{ch:algebra-types}                       -- ch13 (forward)
%
% Difficulty-marker macros (not in alex.sty; defined locally if not yet defined):
\providecommand{\exstar}{$\star$}
\providecommand{\exstarstar}{$\star\star$}
\providecommand{\exstarstarstar}{$\star\star\star$}

\section*{Exercises}
\addcontentsline{toc}{section}{Exercises}

\begin{enumerate}

% -----------------------------------------------------------------------
% WARM-UP  (★ × 3): exercises 1, 2, 3
% -----------------------------------------------------------------------

\item \exstar
\label{ex:ch12-1}
\textbf{Join FPR computation.}
Two FNR-zero approximate relations $\tilde R_1 \subseteq T \times U$
and $\tilde R_2 \subseteq U \times V$ have false-positive rates
$\fprate_1 = 0.05$ and $\fprate_2 = 0.10$. Apply
\Cref{thm:join-rates} under cross-relation independence to compute
$\fprate_{\bowtie}$ for the natural join
$\tilde R_1 \bowtie \tilde R_2$, and report the value to three
significant digits: $\fprate_{\bowtie} = \fprate_1 \cdot \fprate_2
= 0.05 \cdot 0.10 = 0.005$. State $\fnrate_{\bowtie}$ as well: zero,
by the FNR-zero hypothesis on both inputs.

\item \exstar
\label{ex:ch12-2}
\textbf{Projection FPR amplification.}
A relation $\tilde R \subseteq T \times U$ has $\fprate = 0.01$ at
the per-pair level, and the second-attribute domain $U$ has
cardinality $|U| = 100$. From \Cref{sec:relation-composition}, the
projection $\hat t \in \pi_T(\tilde R)$ on a non-projecting $t$
(no $u$ with $(t, u) \in R$) reports $\top$ if any of the $|U|$
independent membership tests reports $\top$. Compute the projection
FPR
\[
  \fprate_{\pi} = 1 - (1 - 0.01)^{100} \approx 1 - 0.366 \approx 0.634
\]
and report the value. State in one sentence the moral: existential
quantification over a moderate-size co-domain inflates a $1\%$
per-pair rate into a $63\%$ per-projection rate.

\item \exstar
\label{ex:ch12-3}
\textbf{Selection composition.}
A relation $\tilde R$ has $\fprate_R = 0.02$ and $\fnrate_R = 0$;
the selection predicate $P$ is exact (deterministic, no error).
From \Cref{sec:relation-composition}, state the output rates of
$\sigma_P(\tilde R)$: the exact predicate passes the rates through
unchanged, so $\fprate_\sigma = 0.02$ and $\fnrate_\sigma = 0$.
In one sentence, contrast with the approximate-predicate case of
exercise~8, where the second factor contributes its own rate by
the chapter~3 AND rule.

% -----------------------------------------------------------------------
% MEDIUM  (★★ × 5): exercises 4, 5, 6, 7, 8
% -----------------------------------------------------------------------

\item \exstarstar
\label{ex:ch12-4}
\textbf{Notebook-tied: MinHash similarity join.}
Open the chapter~12 companion notebook
\texttt{python/notebooks/ch12-approximate-relations.ipynb}.
\begin{enumerate}
  \item[(a)] Generate two synthetic relations $R_1, R_2$ of $1000$
    rows each, each row a token set of size $\approx 10$ drawn so
    that a controlled fraction of cross-relation pairs has true
    Jaccard $J \ge 0.5$ and the rest has $J < 0.5$.
  \item[(b)] Compute MinHash signatures of length $k = 200$ for
    every row, then run the threshold join of
    \Cref{sec:approximate-joins} at $\theta = 0.5$, both with and
    without the LSH band pre-filter at $(b, r) = (20, 10)$.
  \item[(c)] Report empirical $\hat\fprate$ and $\hat\fnrate$ for
    each variant, and compare to the per-pair Hoeffding bound of
    \Cref{prop:minhash-error} and the S-curve survival probability
    $1 - (1 - J^r)^b$ of \Cref{sec:approximate-joins}. In one
    sentence, identify which rate (FPR or FNR) is dominated by the
    LSH layer at this $(b, r)$ setting.
\end{enumerate}

\item \exstarstar
\label{ex:ch12-5}
\textbf{De Morgan for relations.}
The chapter~6 rate algebra (\Cref{prop:union-fpr-fnr},
\Cref{prop:intersection-fpr-fnr}) for approximate sets carries
over to approximate relations on the pair universe $T \times U$,
by \Cref{sec:sets-to-relations}. Show that the De Morgan identity
$\complement(\tilde R_1 \cup \tilde R_2) = \complement \tilde R_1
\cap \complement \tilde R_2$ holds at the rate level: derive the
output $(\fprate, \fnrate)$ of each side under the chapter~6
composition rules and verify they agree, when both inputs have
$\fnrate = 0$ and FPRs $\fprate_1, \fprate_2$ and the
complementation rule swaps $\fprate \leftrightarrow \fnrate$. In
one sentence, identify why the rate-level De Morgan is a corollary
of the classical De Morgan plus the chapter~5 rate definitions
(both sides are the same membership predicate, with the same rate
calculation applied to it).

\item \exstarstar
\label{ex:ch12-6}
\textbf{Three-way join.}
Three FNR-zero approximate relations
$\tilde R_1 \subseteq T_1 \times T_2$,
$\tilde R_2 \subseteq T_2 \times T_3$,
$\tilde R_3 \subseteq T_3 \times T_4$
with FPRs $\fprate_1, \fprate_2, \fprate_3$ are composed by
natural join into $\tilde R_1 \bowtie \tilde R_2 \bowtie \tilde R_3$.
\begin{enumerate}
  \item[(a)] Apply \Cref{thm:join-rates} pairwise: the join of
    the first two has $\fnrate = 0$ and FPR $\fprate_1 \fprate_2$;
    the join with the third has $\fnrate = 0$ and FPR
    $\fprate_1 \fprate_2 \fprate_3$.
  \item[(b)] Conclude that for $k$-fold natural join under
    cross-relation independence, $\fprate_{\bowtie}^{(k)} =
    \prod_{i=1}^k \fprate_i$ and $\fnrate_{\bowtie}^{(k)} = 0$.
    State the analog of the AND-of-Booleans calculation of
    chapter~3: $k$-fold join is $k$-fold AND lifted to the pair
    universe, and the FPR product is the chapter~3 product rule.
  \item[(c)] For $\fprate_1 = \fprate_2 = \fprate_3 = 0.05$,
    compute $\fprate_{\bowtie}^{(3)} = 0.05^3 = 1.25 \times 10^{-4}$.
    Three orders of magnitude of false-positive rate are bought
    for the price of two extra joins, the same gain as the
    chapter~6 intersection of three approximate sets.
\end{enumerate}

\item \exstarstar
\label{ex:ch12-7}
\textbf{LSH parameters for a target precision-recall.}
A similarity-join task at threshold $\theta = 0.7$ requires LSH
parameters $(b, r)$ such that pairs with $J = 0.85$ survive the
LSH filter with probability $\ge 0.95$ (recall target) and pairs
with $J = 0.55$ survive with probability $\le 0.05$ (precision
target). The S-curve survival probability is $1 - (1 - J^r)^b$ per
\Cref{sec:approximate-joins}.
\begin{enumerate}
  \item[(a)] Write the two inequalities:
    $1 - (1 - 0.85^r)^b \ge 0.95$ and $1 - (1 - 0.55^r)^b \le 0.05$.
  \item[(b)] By trial: $r = 4, b = 30$ gives $1 - (1 - 0.85^4)^{30}
    = 1 - (1 - 0.522)^{30} = 1 - 0.478^{30} \approx 1 -
    1.3 \times 10^{-10} \approx 1.000$ on the recall side, and
    $1 - (1 - 0.55^4)^{30} = 1 - (1 - 0.0915)^{30} \approx
    1 - 0.0561 \approx 0.944$ on the precision side, which fails
    the $\le 0.05$ target. Try $r = 8, b = 10$: recall
    $1 - (1 - 0.85^8)^{10} = 1 - (1 - 0.272)^{10} \approx 1 -
    0.0428 \approx 0.957$, precision $1 - (1 - 0.55^8)^{10}
    \approx 1 - 0.918 \approx 0.082$. Iterate to converge on
    $(b, r) \approx (10, 8)$ as a working trade-off.
  \item[(c)] Conclude with the MinHash signature length:
    $k = b \cdot r \approx 80$ for the converged parameters,
    which sets the per-pair cost and the per-pair Hoeffding bound
    of \Cref{prop:minhash-error}.
\end{enumerate}

\item \exstarstar
\label{ex:ch12-8}
\textbf{Composing relation with approximate predicate.}
A relation $\tilde R \subseteq T \times U$ has FPR $\fprate_R =
0.01$ and FNR $\fnrate_R = 0$. The selection predicate $\tilde P$
is itself approximate with FPR $\fprate_P = 0.05$ and FNR
$\fnrate_P = 0$, applied independently of $\tilde R$ per the
cross-relation independence assumption of \Cref{thm:join-rates}.
Apply the chapter~3 AND-under-noise rule (the same rule that
proves \Cref{thm:join-rates}) to derive the joint FPR of
$\sigma_{\tilde P}(\tilde R) = \{(t, u) : \tilde R(t, u) = \top
\text{ and } \tilde P(t, u) = \top\}$:
\[
  \fprate_\sigma = \fprate_R \cdot \fprate_P = 0.01 \cdot 0.05 = 5
    \times 10^{-4},
\]
and state the joint FNR (zero, since both inputs are FNR-zero).
In one sentence, identify the structural fact: AND of two FNR-zero
noisy Booleans always multiplies the FPRs.

% -----------------------------------------------------------------------
% CHALLENGING / EXTENSION  (★★★ × 3): exercises 9, 10, 11
% -----------------------------------------------------------------------

\item \exstarstarstar
\label{ex:ch12-9}
\textbf{Optimality of MinHash for similarity sketches.}
Chapter~8 \S8.3 established an entropy lower bound on the bits per
element needed for an approximate-set predicate with FPR $\le
\epsilon$ to compress its input. The MinHash sketch of
\Cref{def:minhash}, applied to similarity estimation rather than
set membership, has its own information-theoretic floor: the bits
needed to distinguish two sets at Jaccard precision $\pm t$ with
confidence $1 - \delta$ are bounded below by $\Omega(\log(1/t) +
\log(1/\delta))$ per coordinate (an entropy argument on the
collision-probability distribution). Sketch why MinHash achieves
this floor up to constants. \emph{Hint}: each MinHash coordinate
is a uniform draw from the universe restricted to the union $A
\cup B$; its collision probability is exactly $J$, and the
$k$-coordinate signature has Shannon information $k \cdot H(J)$
about the underlying $J$, matching the lower bound in the regime
where the rate function $H(J)$ approaches its argument-derivative
of Hoeffding. The chapter~8 framework gives the lower bound; the
MinHash construction saturates it.

\item \exstarstarstar
\label{ex:ch12-10}
\textbf{Approximate equijoin via Bloom intersection.}
The MinHash-LSH pipeline of \Cref{sec:approximate-joins} handles
similarity joins. The equijoin variant matches rows on equality of
a shared attribute, which is the classical SQL inner join. A
Bloom-filter implementation of the equijoin builds a Bloom filter
$B_2$ over the join-attribute values of $R_2$, scans $R_1$, and
emits $(r_1, r_2)$ pairs where $r_1$'s join attribute hits
$B_2$ and a verification step on the candidate confirms equality.
\begin{enumerate}
  \item[(a)] Sketch the construction: $B_2 \leftarrow$
    \textsc{BloomFilter}($\{r_2.\text{key} : r_2 \in R_2\}$),
    then \textbf{for} $r_1 \in R_1$ \textbf{do if} $r_1.\text{key}
    \in B_2$ \textbf{then} scan $R_2$ for matching $r_2$ values
    and emit pairs.
  \item[(b)] Derive the FPR at the candidate-pair level: a pair
    $(r_1, r_2)$ with $r_1.\text{key} \ne r_2.\text{key}$ becomes
    a candidate iff $r_1.\text{key}$ collides in $B_2$ with some
    stored key, with probability equal to the Bloom-filter
    $\fprate$ from chapter~5. Verification eliminates the
    false-positive pairs at the cost of one comparison per
    candidate; the final classifier has $\fprate_{\bowtie} = 0$
    and $\fnrate_{\bowtie} = 0$ exactly (the Bloom filter is a
    candidate pruner, not a final decision).
  \item[(c)] Identify the framework reading: this is the
    approximate-relation $\tilde R \subseteq R_1 \times R_2$ with
    membership predicate ``$r_1.\text{key}$ hits $B_2$,'' which
    is a one-sided FNR-zero approximate relation in the
    \Cref{sec:relation-composition} sense; the verification step
    is the exact selection of exercise~3 applied to remove the
    residual false positives.
\end{enumerate}

\item \exstarstarstar
\label{ex:ch12-11}
\textbf{Approximate relation as $\Bernoulli{T \times U}$.}
Chapter~13 (\Cref{ch:algebra-types}) recasts the approximate
relations of this chapter as instances of the type-algebraic
construct $\Bernoulli{T \times U}$: the universe is the product
type $T \times U$, and the constructor $\Bernoulli{\cdot}$ acts
functorially on it.
\begin{enumerate}
  \item[(a)] State the identification: an approximate relation
    $\tilde R \subseteq T \times U$ with predicate $\tilde R(t, u)
    \in \Bernoulli{\BoolType}$ equals an element of
    $\Bernoulli{T \times U}$ in the sense that
    $\Bernoulli{T \times U}$ is, by the type-algebraic
    construction, the set of $\Bernoulli{\BoolType}$-valued
    functions on $T \times U$.
  \item[(b)] Identify two composition isomorphisms preserved
    under the recasting: (i) the join $\tilde R_1 \bowtie
    \tilde R_2$ becomes the product-type natural transformation
    that pulls $\Bernoulli{T \times U} \otimes \Bernoulli{U
    \times V}$ to $\Bernoulli{T \times U \times V}$ along the
    shared $U$ axis, and (ii) the projection $\pi_T(\tilde R)$
    becomes the functorial image $\pi_T : \Bernoulli{T \times U}
    \to \Bernoulli{T}$ that traces over $U$ with an existential
    quantifier. State which of the chapter~12 rate formulas
    survives the type-algebra recasting unchanged
    (\Cref{thm:join-rates}: yes; projection FPR amplification:
    yes, as a quantifier-trace formula on the type algebra) and
    which becomes a corollary of a more general algebraic
    identity (the chapter~6 De Morgan rule of exercise~5
    generalizes to all complementable Bernoulli types).
  \item[(c)] In one sentence, identify why the recasting matters:
    the chapter~12 derivations are each separate ``lifts'' of
    chapter~5 and chapter~6, but as type-algebra constructions
    they become single functorial statements that apply
    uniformly across $\Bernoulli{T}, \Bernoulli{T \times U},
    \Bernoulli{L\,T}$, and any other type the algebra produces.
\end{enumerate}

\end{enumerate}
